\section{Manus: de 0 a 100 millones de ARR}

La sección sobre Manus es la que abarca más tramo de la conversación. Lo central aquí no es una cifra de crecimiento llamativa, sino cómo un producto de Agent pasa de la demo a ingresos reales. Peak habla del crecimiento de Manus, de la retroalimentación de los usuarios, del umbral de las herramientas de productividad y del cambio por el que los usuarios dejan de verlo como una «herramienta para ser más eficientes» y lo adoptan como un «factor de producción principal».

\begin{figure}[H]
\centering
\includegraphics[width=0.96\textwidth]{manus-arr-flywheel.png}
\caption{Manus de 0 a 100 millones de ARR: el volante de inercia de producto, usuarios, ingresos, retroalimentación e iteración.}
\end{figure}

\begin{knowledgebox}{Lectura de la figura: detrás del ARR hay un ciclo de uso}
La figura va de la usabilidad del producto a los flujos de trabajo reales, y de ahí a los ingresos, la retroalimentación y la iteración. Un producto de Agent no puede crecer solo con una demostración: tiene que generar valor de forma continua en el trabajo real de los usuarios. El ARR es el resultado de este ciclo, no su punto de partida.
\end{knowledgebox}

\subsection{De la eficiencia al factor de producción}

Peak cuenta que, con cierta versión, los usuarios sintieron que Manus había cruzado el umbral de «herramienta de productividad»: antes, el Agent servía para ayudar a ser más eficiente; después, algunos usuarios lo adoptaron de verdad como su aplicación principal de trabajo para generar ingresos. Este es un punto de quiebre muy importante para los productos de Agent. Una herramienta de eficiencia es sustituible; un factor de producción es mucho más difícil de abandonar.

\begin{importantbox}{El umbral de producto de un Agent}
Un Agent solo entra de verdad en el sistema de productividad cuando pasa de «ayudarme a ahorrar tiempo» a «ayudarme a generar ingresos o a realizar el trabajo central».
\end{importantbox}

\subsection{Cómo la retroalimentación de los usuarios da forma al producto}

En la ronda de preguntas rápidas al final de la entrevista, Peak dice que el hecho clave en este momento es que el siguiente avance de la IA necesita las sugerencias de los usuarios. Esta afirmación es coherente con la lógica de crecimiento de Manus. Un Agent es un producto que interactúa con entornos reales; la retroalimentación de los usuarios no es información de posventa, sino una entrada para el modelo, las herramientas y los límites del producto.

\begin{knowledgebox}{De las sugerencias de los usuarios al aprendizaje del producto}
Las sugerencias de los usuarios pueden revelar flujos de trabajo reales, puntos de falla, necesidades de permisos, necesidades de pago y bloqueos del entorno. Cuanto más cerca está un producto de Agent de las tareas reales, más depende de la retroalimentación de los usuarios para definir el siguiente paso.
\end{knowledgebox}

\subsection{Cuatro pruebas detrás del crecimiento del ARR}

Pasar de 0 a 100 millones de dólares de ARR suena a una cifra comercial, pero para una empresa de Agents se parece más a la suma de cuatro pruebas de producto. Primera, si los usuarios entienden de verdad lo que este producto puede hacer; segunda, si están dispuestos a incorporarlo a su flujo de trabajo principal; tercera, si están dispuestos a pagar por un valor continuo; cuarta, si el producto puede obtener retroalimentación de mayor calidad a partir del uso.

\begin{longtable}{p{0.20\textwidth}p{0.34\textwidth}p{0.37\textwidth}}
\toprule
Prueba & Qué debe demostrar & Qué significa para el Agent \\
\midrule
Comprensible & El usuario puede decir en una frase qué problema resuelve Manus & El producto queda establecido en la mente del usuario; no es solo una demostración técnica. \\
Confiable & El usuario está dispuesto a encargarle a Manus su trabajo central & Pasa de herramienta de eficiencia a sistema de productividad. \\
Pagable & El usuario está dispuesto a pagar de forma continua, no solo a probarlo una vez por novedad & Solo así el ARR tiene efecto compuesto y es predecible. \\
Capaz de aprender & A más uso, más retroalimentación, y el producto entiende mejor las tareas reales & La retroalimentación de los usuarios, a su vez, mejora el modelo y las herramientas. \\
\bottomrule
\end{longtable}

\begin{importantbox}{El ARR es el resultado, no la causa}
El ARR de una empresa de Agents proviene de que «los usuarios le confían su trabajo real al producto». Sin un flujo de trabajo confiable, el crecimiento de los ingresos puede ser fácilmente solo un entusiasmo pasajero; con un flujo de trabajo confiable, el crecimiento puede convertirse en un volante de inercia.
\end{importantbox}

\subsection{Resumen de la sección}

La historia de Manus de 0 a 100 millones de ARR es, en esencia, la de un Agent que pasa de demostración a sistema de productividad. El crecimiento real proviene de que los usuarios estén dispuestos a confiarle su trabajo central.
